% Bagian: computability
% Bab: computability-theory
% Subbagian: ce-closed-cup-cap

\documentclass[../../../include/open-logic-section]{subfiles}

\begin{document}

\olfileid[id]{cmp}{thy}{clo}

\olsection[Gabungan dan Irisan Himpunan C.E.]{Himpunan Terenumerasi Secara
  Komputabel Tertutup terhadap Gabungan dan Irisan}

Teorema berikut memberikan beberapa sifat ketertutupan bagi himpunan-himpunan
yang terenumerasi secara komputabel.

\begin{thm}
Andaikan $A$ dan $B$ terenumerasi secara komputabel. Maka demikian pula
$A \cap B$ dan $A \cup B$.
\end{thm}

\begin{proof}
\olref[eqc]{thm:ce-equiv} memungkinkan kita menggunakan berbagai pencirian
himpunan terenumerasi secara komputabel. Sebagai ilustrasi, kita akan
memberikan beberapa bukti yang berbeda.

Untuk bukti pertama, andaikan $A$ dienumerasi oleh fungsi dapat dihitung~$f$,
dan $B$ dienumerasi oleh fungsi dapat dihitung~$g$. Misalkan
\begin{align*}
h(x) & = \umin{y}{(f(y) = x \lor g(y) = x)} \text{ dan}\\
j(x) & = \umin{y}{(f((y)_0) = x \land g((y)_1) = x)}.
\end{align*}
Maka $A \cup B$ merupakan domain $h$, sedangkan $A \cap B$ merupakan
domain~$j$.

\begin{explain}
Dalam istilah komputasional, yang terjadi adalah sebagai berikut: jika
diberikan prosedur yang mengenumerasi $A$ dan $B$, kita dapat memutuskan-sebagian
apakah suatu unsur $x$ berada dalam $A \cup B$ dengan mencari $x$ dalam salah
satu enumerasi; dan kita dapat memutuskan-sebagian apakah suatu unsur $x$
berada dalam $A \cap B$ dengan mencari $x$ dalam kedua enumerasi secara
bersamaan.
\end{explain}

Untuk bukti kedua, andaikan lagi bahwa $A$ dienumerasi oleh~$f$ dan $B$
dienumerasi oleh~$g$. Misalkan
\[
k(x) = \begin{cases}
f(x/2) & \text{jika $x$ genap} \\
g((x-1)/2) & \text{jika $x$ ganjil.}
\end{cases}
\]
Maka $k$ mengenumerasi $A \cup B$; gagasannya ialah $k$ sekadar berselang-seling
antara enumerasi yang diberikan oleh $f$ dan~$g$. Mengenumerasi $A \cap B$
lebih rumit. Jika $A \cap B$ kosong, himpunan itu secara trivial terenumerasi
secara komputabel. Jika tidak, misalkan $c$ sembarang unsur $A \cap B$, dan
definisikan $l$ dengan
\[
l(x) = \begin{cases}
f((x)_0) & \text{jika $f((x)_0) = g((x)_1)$} \\
c & \text{selain itu.}
\end{cases}
\]
Dalam istilah komputasional, $l$ menelusuri pasangan unsur dalam enumerasi
$f$ dan $g$, lalu mengeluarkan setiap kecocokan yang ditemukan; jika tidak
menemukan kecocokan, $l$ mengulur waktu dengan mengeluarkan $c$.

Untuk bukti terakhir, andaikan $A$ merupakan \emph{domain} fungsi parsial
$m(x)$ dan $B$ merupakan domain fungsi parsial $n(x)$. Maka $A \cap B$
merupakan domain fungsi parsial $m(x) + n(x)$.

\begin{explain}
Dalam istilah komputasional, jika $A$ merupakan himpunan nilai tempat $m$
berhenti dan $B$ merupakan himpunan nilai tempat $n$ berhenti, maka $A \cap B$
adalah himpunan nilai tempat kedua prosedur berhenti.
\end{explain}

Menyatakan $A \cup B$ sebagai himpunan nilai penghentian lebih sulit karena
kita harus menyimulasikan $m$ dan $n$ secara paralel. Misalkan $d$ suatu indeks
bagi $m$ dan $e$ suatu indeks bagi $n$; dengan kata lain, $m = \cfind{d}$ dan
$n = \cfind{e}$. Maka $A \cup B$ merupakan domain fungsi
\[
p(x) = \umin{y}{(T(d,x,y) \lor T(e,x,y))}.
\]
\begin{explain}
Dalam istilah komputasional, pada masukan $x$, $p$ mencari komputasi berhenti
bagi $m$ atau komputasi berhenti bagi $n$, lalu berhenti jika menemukan salah
satunya.
\end{explain}
\end{proof}

\end{document}

